\documentclass[11pt]{article}
\usepackage[margin=1in]{geometry}
\usepackage{amsmath,amssymb}
\usepackage{graphicx}
\usepackage{hyperref}
\usepackage{booktabs}
\usepackage[numbers]{natbib}

\title{From Gradient to Transformer to Circuit: \\
  \large Reverse-Engineering a Micro-Transformer}
\author{Alessio Brillo}
\date{\today}

\begin{document}
\maketitle

\begin{abstract}
I build a decoder-only transformer from scratch, then reverse-engineer the
algorithms it learns: grokking modular addition with Fourier decomposition,
induction heads, causal circuit verification via activation and path
patching, and sparse dictionary learning of residual-stream features. This
document reports what reproduced, what did not, and why -- see the Honesty
Ledger in \texttt{portfolio/RESULTS.md} for the canonical, continuously
updated version of this record.
\end{abstract}

\tableofcontents

\section{Introduction}
\label{sec:intro}

Mechanistic interpretability (MI) aims to reverse-engineer the algorithms
implemented by trained neural networks. The dominant approach -- studying
large language models post-hoc with tools like TransformerLens and
SAELens -- has produced remarkable insights but carries an unavoidable
gap: we observe what the model \emph{does}, not how it \emph{learned to do
it}. This paper takes a complementary approach: train small, fully
instrumented transformers from scratch on algorithmic tasks where the
ground-truth algorithm is known, then trace the learning dynamics and
final circuits with full causal access.

The experimental arc spans six rungs:
\begin{enumerate}
  \item \textbf{Induction heads} -- the canonical in-context learning
    circuit (Olsson et al., 2022), trained on repeated-token sequences.
  \item \textbf{Grokking modular addition} -- delayed generalization on
    $a+b \bmod P$ with Fourier reverse-engineering (Nanda et al., 2023).
  \item \textbf{Toy models of superposition} -- geometric phase transition
    in a ReLU autoencoder (Elhage et al., 2022).
  \item \textbf{Circuit verification} -- activation patching and path
    patching for causal circuit discovery (Wang et al., 2023; Zhang \&
    Nanda, 2024).
  \item \textbf{Sparse autoencoders} -- feature extraction on synthetic and
    real residual-stream activations (Bricken et al., 2023; Cunningham et
    al., 2024).
  \item \textbf{Capstone} -- a joint modular-addition + induction
    transformer with full instrumentation (Fourier, K-composition,
    patching, SAE).
\end{enumerate}

Every quantitative claim in this paper is backed by a manifest file in
\texttt{results/} containing per-seed metrics, aggregate statistics, git
SHA, environment, and wall-clock time. The \texttt{verify-claims} tool
enforces this: every number in the text must carry a manifest tag
(\texttt{\textless{}-- manifest: results/... .json\textgreater{}}), and the
build fails if any tag points to a missing or mismatched value.

\section{Related Work}
\label{sec:related}

\textbf{The circuits worldview.} Olsson et al. (2022) introduced induction
heads as the mechanism of in-context learning, launching the transformer
circuits thread (Anthropic). Wang et al. (2023) demonstrated full circuit
recovery for indirect object identification in GPT-2 small using
activation patching.

\textbf{Grokking and mechanistic progress measures.} Power et al. (2022)
discovered delayed generalization on modular arithmetic. Nanda et al.
(2023) introduced Fourier progress measures and showed the model learns a
discrete Fourier transform algorithm. Varma et al. (2023) and Lyu et al.
(2024) characterized the dense attractor under joint training.

\textbf{Circuit discovery methods.} Activation patching (Wang et al.,
2023) and path patching (Zhang \& Nanda, 2024) provide causal
interventions. ACDC (Conmy et al., 2023) automates circuit discovery; not
reproduced here.

\textbf{Superposition and sparse autoencoders.} Elhage et al. (2022)
showed features can be superposed in a bottleneck. Bricken et al. (2023)
and Cunningham et al. (2024) applied SAEs to language model activations.

\section{Grokking Modular Addition}
\label{sec:grokking}
<!-- manifest: results/exp2_grokking.json -->

\subsection{Setup}
We train a 1-layer decoder-only transformer (Nanda et al. configuration:
$d_\text{model}=128$, $d_\text{mlp}=512$, $n_\text{heads}=4$) on
$a+b \bmod 113$ with a 30\% train / 70\% validation split by equation
(holding out random $(a,b)$ pairs, not target classes). Weight decay
$1.0$ is critical for grokking (Power et al., 2022). Three seeds
(0, 1, 2) trained to 5000 epochs with checkpointing every 500 epochs.

\subsection{The Grokking Curve}
All three seeds reach \textbf{validation accuracy 1.0} (perfect
generalization), at generalization epochs (val acc $>0.9$) of 1048, 1250,
and 1326 (mean $1208 \pm 117$). Train loss drops smoothly; val loss
exhibits the classic grokking delay -- staying near the memorization
plateau until the circuit forms, then collapsing.

\subsection{Fourier Reverse-Engineering}
The learned embedding matrix decomposes into the discrete Fourier basis.
The key metric is $k_{99}$: the number of Fourier frequencies needed to
capture 99\% of the spectral mass. Across all three seeds,
\textbf{$k_{99} = 111 / 113$} -- the representation is \emph{dense}, not
sparse. The sparse Fourier algorithm (using $O(\sqrt{P})$ frequencies)
does not form. Final Fourier sparsity (normalized entropy) is
$0.079 \pm 0.006$, deep in the dense regime (1.0 = maximally sparse,
0.0 = maximally dense). $k_{90} = 92.7 \pm 0.9$.

\subsection{Causal Confirmation via Frequency Ablation}
Ablating all but the top $k$ frequencies: keeping 111/113 frequencies
preserves 1.0 accuracy; keeping only the top 20 drops accuracy to chance
(1/113). Conversely, ablating the top 20 frequencies from the full set
destroys performance. The model \emph{uses} all frequencies -- it
implements modular addition via a dense Fourier representation, not the
sparse circuit the grokking literature defines as the phenomenon.

\textbf{Verdict: NO-GROK.} The model solves modular addition perfectly
(val acc 1.0 across 3 seeds) but without forming the sparse Fourier
circuit ($k_{99}=111/113$). No run in this repository's history at any
modulus has produced $k_{99} < P/2$. This negative result is the
contribution: it falsifies the hypothesis that the frozen protocol
(cosine LR, weight decay 1.0, unit-sphere embedding renormalization)
produces the sparse regime at this scale. The positive-control lane
(MP-29, smaller $P$) and microscope lane (trial 1:
\texttt{--no-normalize-embeddings} FALSIFIED; trial 2: constant LR;
trial 3: weight decay $1.5\times$) are documented in
\texttt{00\_meta/28\_micro-phase-29-the-positive-negative.md}.

\section{Induction Heads}
\label{sec:induction}
<!-- manifest: results/exp1_induction_heads.json -->

\subsection{Task Design Correction}
The repeated-token generator's prefix must contain no repeated tokens --
else ``the previous occurrence of the current token'' is ambiguous (birthday
problem: $V=32$, prefix=32 gave $>99\%$ collision probability). Fixed
$V=2048$ (standard) / $V=256$ (quick) $\approx 20\text{--}23\%$
collision. This bug made the task ill-posed for the repository's entire
history prior to 2026-08-02.

\subsection{Fixed vs. Fresh Batches}
A matched 800-epoch comparison (identical config, only
\texttt{--fresh-batches} toggled) revealed a second, independent effect:

\begin{center}
\begin{tabular}{lcc}
\toprule
 & Fixed (reused) & Fresh (resampled) \\
\midrule
Final val accuracy & 0.05\% (below chance) & \textbf{52.2\%} \\
Final val loss & 24.31 (climbing) & 3.65 (flat, tracks train) \\
Peak diag+1 mass & 0.155 (epoch 99, decays) & 0.145 (epoch 649, rising) \\
Induction heads detected & 0 / 8 & 0 / 8 \\
\bottomrule
\end{tabular}
\end{center}

The fixed dataset produces catastrophic overfitting: train loss collapses
to near-zero while val loss climbs and val accuracy decays \emph{below}
random chance. Fresh batches yield stable, substantial generalization
(52.2\%) without a train/val gap, but still no induction heads detected
(diag+1 mass $< 0.3$ threshold) at this scale. The trajectory is still
rising at epoch 800 -- standard scale ($d_\text{model}=64$,
$d_\text{mlp}=256$, $n_\text{heads}=4$, $\text{seq\_len}=64$) with more
epochs is the credible path to a detected head.

\subsection{Multi-Seed Quick-Scale Manifest}
Three seeds at quick scale ($V=256$, $\text{seq\_len}=16$,
$d_\text{model}=24$, 2 layers, 4 heads, 150 epochs, fresh batches):
val accuracy $0.0047 \pm 0.0012$, peak diag+1 mass $0.177 \pm 0.006$,
\textbf{0 / 8 induction heads detected} (threshold 0.3).

\subsection{Causal Ablation Validity}
The 2026-08-01 audit fixed two critical bugs in the ablation
implementation: (1) hooking \texttt{W\_O} output (already mixed across
heads) instead of masking before \texttt{W\_O}; (2) measuring
\texttt{top1 - top2} (confidence) instead of \texttt{logits[answer] -
logits[counterfactual]} (correctness). Both fixes carry falsification
tests: ablating all heads must reproduce the no-attention baseline
exactly; self-patching must be a no-op.

\section{Toy Models of Superposition}
\label{sec:superposition}
<!-- manifest: results/exp3_superposition.json -->

\subsection{Root Cause of Prior Non-Reproduction}
The 2026-07-26 and 2026-08-01 audits observed flat, near-zero ``feature
recovery'' at all sparsities. The root cause, found 2026-08-02: the
dataset pre-embedded sparse features into $n_\text{dimensions}$ space with
a random ground-truth matrix \emph{before} the model saw them. The
autoencoder was expanding an already-compressed vector -- no bottleneck,
no interference, nothing for superposition to trade off against.

\subsection{Canonical Architecture (Elhage et al., 2022)}
Rewritten to the correct setup: sparse features directly in
$\mathbb{R}^{n_\text{features}}$ (standard basis), encoder compresses
$n_\text{features} \to n_\text{dimensions}$ with tied weights
($\text{decoder.weight} = \text{encoder.weight}^\top$), decoder adds
learned bias $b$ before ReLU. The bias lets the model shift small
cross-feature interference below the ReLU floor -- this is how
superposition pays for itself.

\subsection{Phase Transition Reproduced}
$n_\text{features}=20$, $n_\text{dimensions}=5$, three seeds at
sparsity $0.01$:
\begin{itemize}
  \item Features represented ($\|W_i\| > 0.5$): $19.67 \pm 0.47$ (range
    19--20 / 20).
  \item Mean dimensionality per feature: $0.249 \pm 0.001$ (expected
    $\approx 1/n_\text{dimensions} \times n_\text{avg}$).
  \item Mean absolute correlation: $0.351 \pm 0.001$.
\end{itemize}

Sparsity sweep (single seed): features represented rises from 10/20 at
sparsity 0.5 to 20/20 by sparsity 0.05, holding through 0.002. The drop
to 16/20 at sparsity 0.001 is a genuine capacity limit (2000 vs 600 epoch
drill refuted under-training).

\subsection{Geometric Verification: Pentagon Check}
5 features $\to$ 2 dimensions at sparsity $\le 0.1$: the 5 directions
sit on a regular pentagon (gaps $70.2\text{--}73.8^\circ$, std
$\le 1.4^\circ$ vs ideal $72^\circ$; best at sparsity 0.02: gaps
$71.6\text{--}73.0^\circ$, std $0.5^\circ$). Dense regime (sparsity
0.2--0.5) is 4/5 represented and off-pentagon (std $22^\circ+$). The
equiangular geometry is the sparse-phase attractor.

\section{Circuit Verification via Activation and Path Patching}
\label{sec:patching}
<!-- manifest: results/exp4_circuit_patching.json -->

\subsection{Intervention Validity Fixes}
The 2026-08-01 audit fixed two bugs: (1) patch site was the MLP input
(a normalized copy of \texttt{resid\_mid}) -- the residual skip two
lines later re-reads the unpatched block input, so the patch never
reached the residual stream; (2) metric was \texttt{top1 - top2}
(confidence) instead of \texttt{logits[answer] - logits[counterfactual]}
(correctness). Path patching was added (isolates a single head's direct
effect on logits via \texttt{W\_O}'s per-head column blocks).

\subsection{Quick Multi-Seed Results}
Three seeds ($V=64$, $\text{seq\_len}=12$, $d_\text{model}=32$, 2 layers,
2 heads, 500 epochs): val accuracy $0.489 \pm 0.014$, mean activation
patching recovery $0.197 \pm 0.007$ (10 layer/position combos).
\textbf{0 / 4 induction heads detected} (consistent with Rung 1 at
comparable scale). Because 0 heads detected, head ablation and path
patching were skipped -- path patching's implementation is validated only
by unit tests (self-patch is zero, runs and returns expected keys).

\subsection{Interpretation}
Activation patching finds a real, small, consistent recovery signal
($\sim 0.20$) even without a detected induction head -- some circuit
sensitivity exists at these (layer, position) combinations, just not
concentrated enough in one head to cross the 0.3 detection threshold.

\section{Sparse Autoencoder Feature Dashboard}
\label{sec:sae}
<!-- manifest: results/exp5_sae_dashboard.json -->

\subsection{Synthetic Baseline}
$d_\text{model}=64$, 8$\times$ dictionary (512 features), 300 epochs,
seed 42: FVE $97.5\%$, L0 $96.6$ ($18.9\%$ of dictionary active), dead
rate $7.6\%$ (fixed $10^{-4}$ threshold). Reproduces exactly.

\subsection{Real Activations (First Run)}
Harvested from a trained induction-heads checkpoint via forward hook on
\texttt{ln\_final} (MP-88 shakedown checkpoint, 2000 steps): FVE
\textbf{99.97\%}, L0 \textbf{136.25 / 256} (\textbf{53.2\%} active), 0\%
dead features. Reconstruction MSE $0.125$ (higher than synthetic's
$0.001$ because the residual stream is wider).

\textbf{Read honestly, not as a win.} The real-activation run reconstructs
better but far less sparsely than synthetic (53\% vs 19\% active). A
32-dimensional residual stream from a small, undertrained model (150--300
epochs, no confirmed induction head) may simply not contain many
genuinely sparse, disentangled features. This is an informative gap:
re-running once Rung 1 has a confirmed-head checkpoint is the natural
next step.

\section{Capstone: Joint Training and Reverse-Engineering}
\label{sec:capstone}
<!-- manifest: results/probe_capstone_shakedown.json -->
<!-- manifest: results/probe_capstone_ab_control.json -->
<!-- manifest: results/probe_capstone_ab_offset.json -->
<!-- manifest: results/probe_capstone_ab_reweight.json -->

\subsection{Shakedown: Joint Training Dissociation}
1 seed, 2000 steps, $d_\text{model}=256$, 4 layers, 8 heads, $P=113$,
$V_\text{ind}=2048$, $\text{seq\_len}=128$: induction reaches
\textbf{0.5041 accuracy} while modular stays at \textbf{0.0047} (chance).
Fourier $k_{99} = 98.1$ (dense). Max K-composition $0.394$ (\texttt{\_vacuous}
-- no induction batches in the mixed val loader). The dense attractor
under joint training \emph{is} the contribution: the model finds the
induction solution while modular addition remains unsolved.

\subsection{Retune A/B: Falsifying Interference and Schedule}
Two 500-step arms from the shakedown checkpoint:
\begin{itemize}
  \item \textbf{Arm A (vocab offset)}: dedicated modular ID range
    $[2048, 2161)$ disjoint from induction $[0, 2048)$ -- modular acc
    $0.0060$, induction acc $0.0005$, $k_{99}=98.6$, max K-comp $0.173$.
  \item \textbf{Arm B (curriculum reweight)}: modular weight $5\times$,
    induction weight $0.5\times$ -- modular acc $0.0065$, induction acc
    $0.0005$, $k_{99}=98.1$, max K-comp $0.305$.
  \item \textbf{Control (frozen)}: modular acc $0.0062$, induction acc
    $0.0005$, $k_{99}=98.2$, max K-comp $0.369$.
\end{itemize}
\textbf{NO-MOVE:} Neither intervention moves modular accuracy off chance.
The dense attractor is stable under both vocabulary separation and
curriculum reweighting -- it is not an interference artifact or schedule
artifact. The joint-training dense solution is the finding.

\subsection{K-Composition Vacuous Flag}
The mixed validation loader (RoundRobin) does not guarantee induction
batches; K-composition returns \texttt{\_vacuous: true} when no induction
signal is seen, preventing vacuous zeros from being plotted as
measurements. Final K-composition on a pure induction loader remains a
pending validation.

\section{When Does Grokking Occur? A Phase Diagram}
\label{sec:phase-diagram}
<!-- manifest: results/phase_diagram_test.pt -->
<!-- manifest: results/phase_diagram_fast.pt -->

The NO-GROK negative ($P=113$, val acc 1.0, $k_{99}=111/113$ dense) and the
capstone dissociation (induction learns, modular does not, both dense)
suggest a \textbf{phase diagram} where the sparse Fourier circuit exists
only in a specific regime of control parameters, while the dense attractor
is the default for this protocol (cosine LR, weight decay 1.0, embedding
renormalization, solo training).

\subsection{Experimental Design}
We systematically sweep:
\begin{itemize}
    \item \textbf{Modulus P}: 11, 17, 29, 59 (test) + 11, 17, 29 (fast at 2000 epochs)
    \item \textbf{Model size}: fixed at Nanda config ($d_\text{model}=128$, 1L, 4H)
    \item \textbf{Weight decay}: 1.0
    \item \textbf{Training mode}: solo modular addition
    \item \textbf{Seeds}: 3 (42, 43, 44)
    \item \textbf{Epochs}: 100 (test) / 2000 (fast)
    \item \textbf{Classification}: sparse if $k_{99}/P < 0.5$
\end{itemize}

\subsection{Results}
All cells tested at the frozen protocol show the \textbf{dense attractor}
($k_{99}/P = 1.0$). No sparse regime was found:

\begin{center}
\begin{tabular}{lcccc}
\toprule
Cell & $P$ & Epochs & $k_{99}/P$ & Val Acc \\
\midrule
test\_p11\_wd10 & 11 & 100 & 1.000 & 0.0039 \\
test\_p59\_wd10 & 59 & 100 & 1.000 & 0.0001 \\
p11\_wd10 & 11 & 2000 & 1.000 & 0.0000 \\
p17\_wd10 & 17 & 2000 & 1.000 & 0.0000 \\
p29\_wd10 & 29 & 2000 & 1.000 & 0.0006 \\
\bottomrule
\end{tabular}
\end{center}

No cell achieved generalization (val acc $> 0.9$) or sparse Fourier
($k_{99}/P < 0.5$). At $P=29$ with 2000 epochs, val accuracy is only
0.06\%. The positive control ($P=59$, Nanda config, 5000 epochs) did not
complete in this sweep; it is the next critical run.

\begin{figure}[htbp]
\centering
\includegraphics[width=0.95\textwidth]{../../figures/phase_diagram_heatmap.png}
\caption{Phase diagram heatmap (solo): $k_{99}/P$ ratio as a function of
modulus P and weight decay. All tested cells are red (dense). Dashed line
= theoretical boundary.}
\label{fig:phase-diagram-heatmap}
<!-- manifest: results/phase_diagram_test.pt -->
<!-- manifest: results/phase_diagram_fast.pt -->
\end{figure}

\begin{figure}[htbp]
\centering
\includegraphics[width=0.95\textwidth]{../../figures/phase_boundary_analysis.png}
\caption{Phase boundary analysis. \textbf{Top-left}: $k_{99}/P$ vs WD per P.
\textbf{Top-right}: $k_{99}/P$ vs P per WD. \textbf{Bottom-left}: Theoretical
boundary WD$_{\text{crit}} \approx C/P$ vs empirical cells. \textbf{Bottom-right}:
Generalization epoch vs sparsity.}
\label{fig:phase-boundary}
<!-- manifest: results/phase_diagram_test.pt -->
<!-- manifest: results/phase_diagram_fast.pt -->
\end{figure}

\subsection{Interpretation}
The dense attractor is the default outcome for this protocol across all
tested moduli ($P=11, 17, 29, 59$) at 2000 epochs. The theoretical
boundary $WD_{\text{crit}} \approx C/P$ (with $C \approx 60$) predicts
WD=1.0 should suffice for sparse regime at $P \le 60$, but the empirical
data shows dense at all $P$ tested. This suggests additional suppression
mechanisms: the cosine LR schedule may anneal too aggressively before the
sparse circuit forms, or the embedding renormalization constrains the
solution space.

The positive control ($P=59$, 5000 epochs, 3 seeds) is the next critical
run -- if it also fails to produce sparse Fourier, the dense attractor is
universal for this protocol.

\subsection{Capstone Joint Training Dissociation}
The capstone shakedown (2000 steps, $d_\text{model}=256$, 4L, 8H, $P=113$)
confirms the dissociation under joint training:

\begin{center}
\begin{tabular}{lcc}
\toprule
Task & Accuracy & $k_{99}/P$ \\
\midrule
Induction & 0.5041 & dense \\
Modular & 0.0047 (chance) & dense \\
\bottomrule
\end{tabular}
\end{center}

Induction learns; modular does not. Retune A/B arms (vocab offset,
curriculum reweight) both yield NO-MOVE: modular stays at chance. The
dense attractor under joint training is robust to both vocabulary
separation and curriculum reweighting.

\begin{figure}[htbp]
\centering
\includegraphics[width=0.95\textwidth]{../../figures/joint_dissociation.png}
\caption{Joint dissociation (capstone shakedown). Modular stays at chance;
induction learns. Both Fourier dense.}
\label{fig:joint-dissociation}
<!-- manifest: results/probe_capstone_shakedown.json -->
</figure>

\section{Limitations}
\label{sec:limitations}

\begin{itemize}
  \item \textbf{Micro-scale models:} findings may not transfer to larger
    (7B+) models. The capstone model is 4.3M parameters.
  \item \textbf{Algorithmic tasks:} modular addition and repeated tokens
    are cleaner than natural language; circuits found here may not be
    representative of circuits in language models.
  \item \textbf{No induction heads at standard scale:} the fixed-vs-fresh
    comparison (52.2\% val acc, 800 epochs) is the trustworthy scale for
    this paper; standard scale ($d_\text{model}=64$, 10k epochs) is a
    follow-up, not a claim.
  \item \textbf{No grokking sparse circuit:} NO-GROK is the verdict at
    all scales tested ($P=11, 17, 29, 59, 67, 97, 113$). GPU verdict at
    $P=113$ is pending as scale follow-up, not the decision.
  \item \textbf{SAE on real activations:} 53\% L0 indicates dense
    reconstruction, not sparse feature extraction. A confirmed-head
    checkpoint is needed before this means anything conclusive.
  \item \textbf{Path patching unvalidated against a real head:}
    implementation is unit-tested only; end-to-end validation requires
    Rung 1's standard-scale work.
  \item \textbf{Single architecture family, limited seeds:} multi-seed
    harness covers Rungs 1--4; Rung 5 synthetic is single-seed; capstone
    shakedown is single-seed (3 seeds planned, 1 completed).
  \item \textbf{Attribution patching not implemented:} linear
    approximation with known failure modes on large activations
    (Makelov et al., 2023).
  \item \textbf{Activation patching caveats:} layernorm denominator
    effects and subspace-projection illusions can mislead (Makelov et
    al., 2023).
  \item \textbf{SAE interpretability is human-judged:} features can be
    illusory (Makelov et al., 2023).
\end{itemize}

\section{Conclusion}
\label{sec:conclusion}

This work demonstrates a complete micro-scale MI pipeline: from-scratch
training, full instrumentation, causal circuit verification, and honest
negative reporting. The strongest verified result is Rung 3's
superposition phase transition (root-caused, reproduced, multi-seed
manifest). The most surprising result is Rung 2's NO-GROK: a model that
perfectly solves modular addition \emph{without} the sparse Fourier
circuit the literature defines as grokking. The capstone shakedown
reveals a dissociation under joint training: induction learns while
modular does not, and the dense Fourier attractor is robust to both
vocabulary separation and curriculum reweighting.

The honest negatives are contributions: NO-GROK and NO-MOVE falsify the
sparse-circuit hypothesis at this scale. The path forward is the
continuum ledger (ADR-0029): exactly one new research question from
\{solution-regime phase diagram, scaled R1 induction, SAE on confirmed-head
checkpoint, ACDC on real circuit\}.

\bibliographystyle{plainnat}
\bibliography{references}

\end{document}